\scanpage{165}
\mspage{\deleted{\unsure{2}}39}

Les hypothèses \inserted{\unsure{sur} $W$\unsure{(|C|,}} faites \unsure{tantôt} sur $\unsure{\mathcal M|M}$ \deleted{\unsure{assurent}} servent
à assurer que si je prends $C_0$ petit, \unsure{donc|alors} la
factorisation marche pour $(C,TF)$ automatiquement. Il reste le problème pour
$TC,F$.

Notons que si elle marche, alors $TC=\widetilde{TC}=F^*$,
donc $TC$ est clos à gauche, et a fortiori
satisfait les conditions de stabilité correspondantes
\[
\left\{\begin{aligned}
&\text{a) Stable par cochgt. de base, par composition, par}\\
&\qquad\text{facteurs directs (ces \inserted{deux} derniers sont assurés}\\
&\qquad\text{automatiquement, car $C,W$ sont stables),}\\
&\text{b) Stable par $\unsure{\underrightarrow{\lim}|\underline{\lim}}$ ordinales strictes. \deleted{Ce qui}}
\end{aligned}\right. \tag*{(4.4.3)}
\]
\deleted{\unsure{à cet égard on}} L'hypothèse \inserted{de stabilité} sur $W$ sert \unsure{(}entre autres) à garantir
\deleted{que} la stabilité b) pour $TC$ (en fait, il
suffirait pour cela \inserted{seul} que $W$ soit stable par
$\underrightarrow{\lim}$ ordinales strictes). Il reste la stabilité
de $TC=C\cap W$ par cochgt. de base, qui
est problématique, car $W$ n'est en général
\uline{pas} stable par cochgt. de base. C'est pour
assurer cette stabilité que je me vois
pratiquement contraint ici à imposer
\[
 C_0\subset\mathrm{Cof}_W,\tag*{(4.4.\deleted{\unsure{3}}\unsure{4|3})}
\]
je dis que sous les hypothèses de stabilité
faites sur $W$, cela implique
\[
 C\overset{\mathrm{déf}}{=}\widetilde C_0\subset\mathrm{Cof}_W.\tag*{(4.4.5)}
\]

% archive sheet number: 164
